\begin{frame}
    {Example. Factoring $n=55$}
	In Table \ref{tab:table} we illustrate the classical treatment for some of the most probable results $\tilde{c}/q$ of the quantum subroutine:

	\begin{table}[H]
		\centering
		\begin{tabular}{c c c c}
			\hline
			$\tilde{c}/q$ & $d/r$ & $r$ & Relative frequency  \\ 
			\hline
			0.04993 & 1/20 & 20 & 0.029 \\ 
			0.09998 & 1/10 & 10 & 0.045 \\ 
			0.80005 & 4/5  & 5  & 0.029 \\ 
			0.25000 & 1/4  & 4  & 0.049  \\ 
			\hline
		\end{tabular}
		\caption{Simulation results of $\tilde{c}/q$ with their corresponding relative frequency, the resulting $d/r$ from continued fractions expansion and the estimation for $r$.}
				\label{tab:table}
	\end{table}
	
	\begin{itemize}
		\item One source of failure is that $d$ and $r$ might share common factors that simplify. This occurs precisely when the results $r=10,5,4$ are obtained.
	\end{itemize}
\end{frame}